\documentclass[11pt,a4paper]{article}

% --- Kodowanie i język (polskie znaki) ---
\usepackage[T1]{fontenc}
\usepackage[utf8]{inputenc}
\usepackage[polish]{babel}
\usepackage{lmodern}
\usepackage{microtype}

% --- Układ strony ---
\usepackage[a4paper, left=2.5cm, right=2.5cm,
            top=2.5cm, bottom=2.5cm]{geometry}
\usepackage{setspace}
\setstretch{1.08}
\setlength{\parindent}{1.25cm}
\setlength{\parskip}{2pt}

% --- Matematyka, tabele i grafiki ---
\usepackage{amsmath,amssymb}
\usepackage{graphicx}
\usepackage{float}
\usepackage{booktabs}
\usepackage{array}
\usepackage{tabularx}
\usepackage{longtable}
\usepackage{caption}
\usepackage{subcaption}
\usepackage{placeins}

% --- Listy i kod ---
\usepackage{enumitem}
\setlist[itemize]{topsep=3pt,itemsep=2pt,parsep=0pt}
\usepackage{xcolor}
\usepackage{listings}

\lstdefinestyle{sqlstyle}{
    language=SQL,
    basicstyle=\ttfamily\small,
    keywordstyle=\color{blue!65!black}\bfseries,
    commentstyle=\color{gray!80!black}\itshape,
    stringstyle=\color{red!55!black},
    numbers=left,
    numberstyle=\tiny\color{gray},
    stepnumber=1,
    numbersep=8pt,
    showstringspaces=false,
    breaklines=true,
    frame=single,
    rulecolor=\color{black!25},
    tabsize=4,
    captionpos=b
}
\lstset{style=sqlstyle}

% --- Nagłówki, linki i odsyłacze ---
\usepackage{fancyhdr}
\pagestyle{fancy}
\fancyhf{}
\fancyfoot[C]{\thepage}
\renewcommand{\headrulewidth}{0pt}
\renewcommand{\footrulewidth}{0pt}
\usepackage[hidelinks]{hyperref}
\usepackage[nameinlink,noabbrev]{cleveref}

% --- Ustawienia podpisów ---
\captionsetup{font=small,labelfont=bf}
\renewcommand{\arraystretch}{1.15}

% ============================================================
% DANE PROJEKTU – ZMIEŃ W TYM MIEJSCU
% ============================================================
\newcommand{\tytulprojektu}{BDBT Projekt - Klub Lekkoatletyczny}
\newcommand{\autorzy}{Emilia Wilczek 343170\par Piotr Małowiecki 337127}
\newcommand{\dataoddania}{\today}

% Pomocnicza kolumna tabel o stałej szerokości
\newcolumntype{L}[1]{>{\raggedright\arraybackslash}p{#1}}
\newcolumntype{C}[1]{>{\centering\arraybackslash}p{#1}}

\begin{document}

% ============================================================
% STRONA TYTUŁOWA
% ============================================================
\begin{titlepage}
    \centering

    {\LARGE\bfseries \tytulprojektu\par}
    \vspace{0.6cm}
    {\large \autorzy\par}

    \vspace{1.5cm}
    {\large \dataoddania\par}
    \vspace*{1cm}
\end{titlepage}

% ============================================================
% SPIS TREŚCI
% ============================================================
\tableofcontents
\clearpage

% ============================================================
% 1. WPROWADZENIE
% ============================================================
\section{Zakres i cel projektu}
% Opisz problem, cel projektu, zakres modelowanej dziedziny,
% przyjęte założenia oraz użyte narzędzia.
Celem projektu jest modelowanie oraz implementacja bazy danych w modelu relacyjnym. Projekt składa
się z trzech części: - modelu konceptualnego, modelu logicznego oraz implementacji fizycznej w
silniku bazy danych Oracle XE. Baza danych modelowana jest w ramach zarządzania klubem
lekkoatletycznym, jego sportowcami, ich trenerami, oraz przyrządami do ćwiczeń, informacją o
lokalizacji poszczególnych przyrządów i kiedy są dostępne.

% ============================================================
% 2. DEFINICJA SYSTEMU
% ============================================================
\section{Definicja systemu}

\subsection{Perspektywy użytkowników}
Do głównych użytkowników naszego systemu przypisujemy następujące perspektywy:
\begin{itemize}
    \item \textbf{Administrator} -- Pełny wgląd, modyfikacja list pracowników, stadionów i przyrządów.
    \item \textbf{Trener} -- Wgląd w sportowców których trenuje oraz gdzie i kiedy odbywają się treningi.
    \item \textbf{Sportowiec} -- Wgląd w swoich trenerów, gdzie i kiedy odbywają się treningi, a także zapis na kolejne treningi.
\end{itemize}

\subsection{Funkcjonalności zaimplementowanego systemu}

\subsubsection{Administrator}
\begin{itemize}
    \item Dodawanie/modyfikacja/usuwanie oferowanych ćwiczeń.
    \item Dodawanie/modyfikacja/usuwanie pracowników.
    \item Dodawanie/modyfikacja/usuwanie przyrządów.
    \item Dodawanie/modyfikacja/usuwanie sportowców.
    \item Dodawanie/modyfikacja/usuwanie stadionów.
\end{itemize}

\subsubsection{Trener}
\begin{itemize}
    \item Podgląd swojego grafiku.
    \item Dodawanie/modyfikacja stanu technicznego przyrządu.
\end{itemize}

\subsubsection{Sportowiec}
\begin{itemize}
    \item Podgląd swojego grafiku.
    \item Dodawanie/modyfikacja stanu technicznego przyrządu.
    \item Dodawanie do grafku (zapis na kolejne treningi).
\end{itemize}

\clearpage
% ============================================================
% 3. MODEL KONCEPTUALNY
% ============================================================
\section{Model konceptualny}

\subsection{Definicja zbiorów encji}
Dla każdej encji opisz jej rolę w modelowanej dziedzinie.
\begin{itemize}
    \item \textbf{Klub} -- Główna encja, klub lekkoatletyczny.
    \item \textbf{Pracownik} -- Osoba zatrudniona przez klub.
    \item \textbf{Trener} -- Pracownik który jest trenerem.
    \item \textbf{Cwiczenie} -- Ćwiczenie które jest oferowane przez klub.
    \item \textbf{Stadion} -- Lokacja fizyczna gdzie odbywają się ćwiczenia.
    \item \textbf{Sportowiec} -- Członek klubu, osoba która wykonuje ćwiczenia.
    \item \textbf{Grafik} -- Grafik sesji treningowych.
    \item \textbf{Przyrzad} -- Przyrząd do ćwiczeń.
\end{itemize}

\subsection{Ustalenie związków między encjami i ich typów}
\begin{table}[H]
\centering
\small
\begin{tabularx}{\textwidth}{|L{3.2cm}|C{1.6cm}|X|}
\hline
\textbf{Relacja} & \textbf{Rodzaj} & \textbf{Opis} \\
\hline
Klub -- Pracownik & 1:N & Klub zatrudnia procowników \\
\hline
Pracownik -- Trener & 1:1 & Dziedziczenie: pracownik może być trenerem. Dane trenera są
rozszerzeniem danych pracownika. \\
\hline
Cwiczenie -- Grafik & 1:N & Jakim ćwiczeniem jest dany wpis w grafiku. \\
\hline
Cwiczenie -- Sportowiec & M:N & Ćwiczenia które wykonuje sportowiec. \\
\hline
Cwiczenie -- Stadion & M:N & Ćwiczenia które odbywają się na danych stadionach. \\
\hline
Klub -- Cwiczenie & 1:N & Klub oferuje dane ćwiczenia. \\
\hline
Klub -- Pracownik & 1:N & Klub zatrudnia pracownika. \\
\hline
Klub -- Stadion & 1:N & Klub posiada stadion. \\
\hline
Pracownik -- Grafik & 1:N & Pracownik (który jest trenerem) jest przypisany do danej sesji treningowej w grafiku. \\
\hline
Pracownik -- Sportowiec & M:N & Pracownik (który jest trenerem) trenuje sportowca. \\
\hline
Pracownik -- Stadion & M:N & Pracownik (który jest trenerem) trenuje sportowców na danych stadionach. \\
\hline
Przyrzad -- Grafik & 1:N & Jaki przyrząd jest urzywany w danej sesji treningowej \\
\hline
Sportowiec -- Grafik & 1:N & Zapisy sportowca na sesje treningowe. \\
\hline
Stadion -- Grafik & 1:N & Na którym stadionie odbywa się konkretna sesja trningowa. \\
\hline
Stadion -- Sportowiec & M:N & Sportowiec chodzi na dany stadion. \\
\hline
\end{tabularx}
\caption{Związki między encjami}
\label{tab:zwiazki}
\end{table}

\subsection{Określenie atrybutów i ich dziedzin}
% Skopiuj poniższą tabelę dla każdej encji.
\subsubsection{Klub}
\begin{table}[H]
\centering
\small
\begin{tabularx}{\textwidth}{|L{3cm}|L{2.25cm}|C{2.67cm}|X|}
\hline
\textbf{Atrybut} & \textbf{Typ} & \textbf{Obowiązkowy} & \textbf{Opis} \\
\hline
ID\_klubu & Int & tak & Unikalny identyfikator encji. \\
\hline
Nazwa & Varchar(50) & tak & Nazwa klubu. \\
\hline
Adres & Varchar(100) & tak & Adres siedziby klubu. Pole segmentowe zawierające miasto, ulicę,
numer budynku, numer lokalu oraz adres pocztowy. \\
\hline
REGON & Varchar(14) & nie & Numer REOGON klubu. \\
\hline
Data\_zalozenia & Date & nie & Data założenia klubu. \\
\hline
\end{tabularx}
\caption{Atrybuty encji Klub}
\label{tab:atrybuty-encja-klub}
\end{table}

\subsubsection{Pracownik}
\begin{table}[H]
\centering
\small
\begin{tabularx}{\textwidth}{|L{3cm}|L{2.25cm}|C{2.67cm}|X|}
\hline
\textbf{Atrybut} & \textbf{Typ} & \textbf{Obowiązkowy} & \textbf{Opis} \\
\hline
ID\_pracownika & Int & tak & Unikalny identyfikator encji. \\
\hline
Imie & Varchar(50) & tak & Imię pracownika. \\
\hline
Nazwisko & Varchar(50) & tak & Nazwisko pracownika. \\
\hline
Stanowisko & Varchar(50) & tak & Nazwa stanowiska pracownika. \\
\hline
Data\_zatrudnienia & Date & tak & Data zatrudnienia pracownika. \\
\hline
Nr\_telefonu & Char(9) & nie & Numer telefonu pracownika. \\
\hline
\end{tabularx}
\caption{Atrybuty encji Pracownik}
\label{tab:atrybuty-encja-pracownik}
\end{table}

\subsubsection{Trener}
\begin{table}[H]
\centering
\small
\begin{tabularx}{\textwidth}{|L{3cm}|L{2.25cm}|C{2.67cm}|X|}
\hline
\textbf{Atrybut} & \textbf{Typ} & \textbf{Obowiązkowy} & \textbf{Opis} \\
\hline
Specjalizacja & Varchar(100) & tak & Dysycyplina sportowa w jakiej się specjalizuje. \\
\hline
Nr\_licencji & Char(14) & tak & Numer licencji Polskiego Związku Lekkiej Atletyki. \\
\hline
Waznosc\_licencji & Date & tak & Data ważności licencji. \\
\hline
Staz\_lat & Int & nie & Ilość lat stażu. \\
\hline
\end{tabularx}
\caption{Atrybuty encji Trener}
\label{tab:atrybuty-encja-trener}
\end{table}

\subsubsection{Cwiczenie}
\begin{table}[H]
\centering
\small
\begin{tabularx}{\textwidth}{|L{3cm}|L{2.25cm}|C{2.67cm}|X|}
\hline
\textbf{Atrybut} & \textbf{Typ} & \textbf{Obowiązkowy} & \textbf{Opis} \\
\hline
ID\_cwiczenia & Int & tak & Unikalny identyfikator encji. \\
\hline
Nazwa & Varchar(50) & tak & Nazwa dysycyplina sportowej. \\
\hline
Czas\_trwania\_min & Int & tak & Czas trwania sesji w minutach. \\
\hline
Poziom\_trudnosci & Int & tak & Poziom trudności ćwiczenia. \\
\hline
Liczba\_powtorzen & Int & nie & Liczba powtórzeń w sesji. \\
\hline
Opis & Text & nie & Opis ćwiczenia. \\
\hline
\end{tabularx}
\caption{Atrybuty encji Cwiczenie}
\label{tab:atrybuty-encja-cwiczenie}
\end{table}

\subsubsection{Stadion}
\begin{table}[H]
\centering
\small
\begin{tabularx}{\textwidth}{|L{3cm}|L{2.25cm}|C{2.67cm}|X|}
\hline
\textbf{Atrybut} & \textbf{Typ} & \textbf{Obowiązkowy} & \textbf{Opis} \\
\hline
ID\_stadionu & Int & tak & Unikalny identyfikator encji. \\
\hline
Nazwa & Varchar(50) & tak & Nazwa stadionu. \\
\hline
Adres & Varchar(100) & tak & Adres stadionu. Pole segmentowe zawierające miasto, ulicę,
numer budynku, numer lokalu oraz adres pocztowy. \\
\hline
Typ\_obiektu & Varchar(50) & tak & Typ obiektu. \\
\hline
Nawierzchnia & Varchar(50) & tak & Rodzaj nawierzchni obecny w tym obiekcie. \\
\hline
Godziny\_otwarcia & Varchar(255) & nie & Pole złożone składające się z godzin otwarcia w
poszczególnych dniach tygodnia. \\
\hline
\end{tabularx}
\caption{Atrybuty encji Stadion}
\label{tab:atrybuty-encja-stadion}
\end{table}

\subsubsection{Sportowiec}
\begin{table}[H]
\centering
\small
\begin{tabularx}{\textwidth}{|L{3cm}|L{2.25cm}|C{2.67cm}|X|}
\hline
\textbf{Atrybut} & \textbf{Typ} & \textbf{Obowiązkowy} & \textbf{Opis} \\
\hline
ID\_stadionu & Int & tak & Unikalny identyfikator encji. \\
\hline
Nazwa & Varchar(50) & tak & Nazwa stadionu. \\
\hline
Adres & Varchar(100) & tak & Adres stadionu. Pole segmentowe zawierające miasto, ulicę,
numer budynku, numer lokalu oraz adres pocztowy. \\
\hline
Typ\_obiektu & Varchar(50) & tak & Typ obiektu. \\
\hline
Nawierzchnia & Varchar(50) & tak & Rodzaj nawierzchni obecny w tym obiekcie. \\
\hline
Godziny\_otwarcia & Varchar(255) & nie & Pole złożone składające się z godzin otwarcia w
poszczególnych dniach tygodnia. \\
\hline
\end{tabularx}
\caption{Atrybuty encji Sportowiec}
\label{tab:atrybuty-encja-sportowiec}
\end{table}

\subsubsection{Grafik}
\begin{table}[H]
\centering
\small
\begin{tabularx}{\textwidth}{|L{3cm}|L{2.25cm}|C{2.67cm}|X|}
\hline
\textbf{Atrybut} & \textbf{Typ} & \textbf{Obowiązkowy} & \textbf{Opis} \\
\hline
Czas\_rozpoczecia & Date & tak & Data i godzina rozpoczęcia sesji. \\
\hline
Czas\_zakonczenia & Date & tak & Data i godzina zakończenia sesji. \\
\hline
\end{tabularx}
\caption{Atrybuty encji Grafik}
\label{tab:atrybuty-encja-grafik}
\end{table}

\subsubsection{Przyrzad}
\begin{table}[H]
\centering
\small
\begin{tabularx}{\textwidth}{|L{3cm}|L{2.25cm}|C{2.67cm}|X|}
\hline
\textbf{Atrybut} & \textbf{Typ} & \textbf{Obowiązkowy} & \textbf{Opis} \\
\hline
ID\_przyrzadu & Int & tak & Unikalny identyfikator encji. \\
\hline
Rodzaj & Varchar(50) & tak & Rodzaj przyrządu. \\
\hline
Nazwa & Varchar(50) & tak & Nazwa przyrządu. \\
\hline
Wynajety & Bool & tak & Czy dany przyrząd jest wynajęty. \\
\hline
Stan\_techniczny & Varchar(50) & nie & Określenie stanu technicznego. \\
\hline
Data\_zakupu & Date & nie & Data zakupu przyrządu. \\
\hline
\end{tabularx}
\caption{Atrybuty encji Przyrząd}
\label{tab:atrybuty-encja-przyrzad}
\end{table}

\subsection{Dodatkowe reguły integralnościowe (reguły biznesowe)}
\begin{itemize}
    \item Pracownik może prowadzić treningi wyłącznie wtedy, gdy posiada odpowiednie kwalifikacje.
    \item Pracownik nie może mieć daty zatrudnienia późniejszej niż aktualna data.
    \item Pracownik może pracować na kilku stadionach należących do klubu, ale nie może prowadzić
        dwóch zajęć w tym samym czasie.
    \item Każdy trener musi posiadać ważną licencję trenerską.
    \item Data ważności licencji trenera nie może być wcześniejsza niż data rozpoczęcia prowadzonych
        przez niego zajęć.
    \item Data rozpoczęcia treningu musi być wcześniejsza niż data jego zakończenia.
    \item Sportowiec nie może zapisać się na dwa treningi odbywające się w tym samym czasie.
\end{itemize}

\subsection{Klucze kandydujące i główne (decyzje projektowe)}
\begin{table}[H]
\centering
\small
\begin{tabularx}{\textwidth}{|X|X|X|}
\hline
\textbf{Nazwa encji} & \textbf{Klucz główny} & \textbf{Klucz kandydujący} \\
\hline
Klub & ID\_klubu & NIP, REGEON (jeżeli podany) \\
\hline
Pracownik & ID\_pracownika & -- \\
\hline
Trener & ID\_prcownika & Nr\_licencji \\
\hline
Cwiczenie & ID\_cwiczenia & -- \\
\hline
Stadion & ID\_stadionu & -- \\
\hline
Sportowiec & ID\_Sportowca & -- \\
\hline
Grafik & ID\_wpisu & -- \\
\hline
Przyrzad & ID\_Przyrzadu & -- \\
\hline
\end{tabularx}
\caption{Klucze główne i kandydujące}
\label{tab:klucze}
\end{table}

\subsection{Schemat ER na poziomie konceptualnym}
\begin{figure}[H]
    \centering
    \includegraphics[width=1\textwidth]{raport.t2d/model.pdf}
    \caption{Model konceptualny}
    \label{fig:er-konceptualny}
\end{figure}

\subsection{Problem pułapek szczelinowych i wachlarzowych – analiza i przykłady}
\subsubsection{Pułapka szczelinowa}
Może pojawić się w przypadku relacji pomiędzy encjami Klub, Pracownik i Stadion. Klub zatrudnia
pracowników, którzy pracują na określonych stadionach, a jednocześnie klub posiada stadiony.
Jeżeli nie każdy stadion ma przypisanego pracownika, wyszukiwanie stadionów za pośrednictwem
tabeli Pracownik może zwrócić niekompletne wyniki. Bezpośrednia relacja między klubem a stadionem
pozwala uzyskać pełniejsze informacje.

\subsubsection{Pułapka wachlarzowa}
może występować w przypadku encji Klub, która jest powiązana z encjami Pracownik oraz Ćwiczenie.
Obie relacje mają charakter 1:N. Podczas jednoczesnego łączenia tych tabel może dojść do
zwielokrotnienia wyników zapytania SQL. Przykładowo pracownicy mogą zostać połączeni ze wszystkimi
ćwiczeniami oferowanymi przez dany klub, mimo że nie prowadzą wszystkich tych zajęć. Problem można
rozwiązać przez odpowiednie odwzorowanie relacji między pracownikami a ćwiczeniami, np. za pomocą
tabeli pośredniej.

% ============================================================
% 4. MODEL LOGICZNY
% ============================================================
\section{Model logiczny}

\subsection{Charakterystyka modelu relacyjnego}
[Opisz przejście od modelu konceptualnego do relacyjnego.]

\subsection{Zmiany właściwości niekompatybilnych z modelem relacyjnym}
\begin{itemize}
    \item [Opis zmiany pierwszej.]
    \item [Opis zmiany drugiej.]
    \item [Opis sposobu odwzorowania relacji M:N.]
\end{itemize}

\subsection{Proces normalizacji -- analiza i przykłady}
\subsubsection{Pierwsza postać normalna (1PN)}
[Uzasadnij spełnienie warunków 1PN.]

\subsubsection{Druga postać normalna (2PN)}
[Uzasadnij spełnienie warunków 2PN.]

\subsubsection{Trzecia postać normalna (3PN)}
[Uzasadnij spełnienie warunków 3PN.]

\subsection{Schemat ER na poziomie modelu logicznego}
\IfFileExists{images/er_logiczny.png}{%
\begin{figure}[H]
    \centering
    \includegraphics[width=\textwidth,height=0.68\textheight,keepaspectratio]{images/er_logiczny.png}
    \caption{Schemat logiczny bazy danych}
    \label{fig:er-logiczny}
\end{figure}
}{%
\noindent\textit{[Wstaw schemat ER w pliku graficznym model-logiczny.]}
}

\subsection{Więzy integralności}
\begin{itemize}
    \item \textbf{Integralność encji:} [opis kluczy głównych].
    \item \textbf{Integralność referencyjna:} [opis kluczy obcych].
    \item \textbf{Integralność domenowa:} [opis typów i ograniczeń].
\end{itemize}

\subsection{Proces denormalizacji -- analiza i przykłady}
[Opisz rozważone przypadki denormalizacji oraz uzasadnij decyzje.]

% ============================================================
% 5. FAZA FIZYCZNA
% ============================================================
\section{Faza fizyczna}

\subsection{Projekt transakcji i weryfikacja ich wykonalności}
\begin{table}[H]
\centering
\small
\begin{tabularx}{\textwidth}{|L{5.4cm}|L{2.6cm}|L{2.3cm}|X|}
\hline
\textbf{Funkcjonalność} & \textbf{Uprawnienia} & \textbf{Wykonalność} & \textbf{Potrzebne zasoby} \\
\hline
Dodawanie/modyfikacja/usuwanie oferowanych ćwiczeń & Administrator & Wykonalne & Tabela Cwiczenie \\
\hline
Dodawanie/modyfikacja/usuwanie pracowników & Administrator & Wykonalne & Tabela Pracownik \\
\hline
Dodawanie/usuwanie przyrządów & Administrator & Wykonalne & Tablea Przyrzad \\
\hline
Modyfikacja przyrządów & Administrator, Trener, Sportowiec & Wykonalne & Tabela Przyrzad \\
\hline
Dodawanie/modyfikacja/usuwanie sportowców & Administrator & Wykonalne & Tabela Sportowiec \\
\hline
Dodawanie/modyfikacja/usuwanie stadionów & Administrator & Wykonalne & Tabela Stadion \\
\hline
Podgląd swojego grafiku & Trener, Sportowiec & Wykonalne & Tabela Grafik \\
\hline
Zapis na kolejne treningi & Sportowiec & Wykonalne & Tabela Grafik \\
\hline
\end{tabularx}
\caption{Weryfikacja wykonalności transakcji}
\label{tab:transakcje}
\end{table}

\subsection{Strojenie bazy danych -- dobór indeksów}
\begin{lstlisting}
TODO TODO add indexes for all foreign keys

-- idx_grafik_czas_rozpoczecia: Przyspieszenie wyszukiwania bliskich cwiczen
CREATE INDEX idx_grafik_czas_rozpoczecia ON Grafik(Czas_rozpoczecia);

-- idx_przyrzad_stan_techniczny: Wyszukiwanie uszkodzonych sprzetow itd.
CREATE INDEX idx_przyrzad_stan_techniczny ON Przyrzad(Stan_techniczny);
\end{lstlisting}

\subsection{Skrypt SQL zakładający bazę danych}
% Dłuższy skrypt można dołączyć jako osobny plik:
% \lstinputlisting[caption={Skrypt tworzący bazę danych}]{sql/create_database.sql}
\begin{lstlisting}[caption={Przykładowa definicja tabeli}]
CREATE TABLE Przykladowa_tabela (
    id          NUMBER PRIMARY KEY,
    nazwa       VARCHAR2(100) NOT NULL,
    data_utw    DATE
);
\end{lstlisting}

\subsection{Przykłady zapytań i poleceń SQL}
Każde zapytanie opisz krótko: jaki problem rozwiązuje i jaki wynik zwraca.
\begin{lstlisting}[caption={Przykładowe zapytanie}]
SELECT kolumna, COUNT(*) AS liczba
FROM Nazwa_tabeli
GROUP BY kolumna;
\end{lstlisting}

% ============================================================
% DODATKOWE ELEMENTY – OPCJONALNE
% ============================================================
% Dodaj kolejne sekcje, jeśli wymaga tego temat lub prowadzący.
% Przykładowo:
%
% \section{Testowanie}
% \subsection{Dane testowe}
% \subsection{Scenariusze testowe}
%
% \section{Podsumowanie}
% [Wnioski końcowe, ograniczenia i możliwe kierunki rozwoju.]

% ============================================================
% BIBLIOGRAFIA – JEŚLI WYMAGANA
% ============================================================
% \begin{thebibliography}{9}
% \bibitem{oracle} Oracle, dokumentacja systemu Oracle Database,
% \url{https://docs.oracle.com/}.
% \end{thebibliography}

\end{document}
